% Part: sets-functions-relations
% Chapter: sets
% Section: important-sets

\documentclass[../../../include/open-logic-section]{subfiles}

\begin{document}

\olfileid[fa-IR]{sfr}{set}{imp}
\olsection{برخی مجموعه‌های مهم}

\begin{ex}
بیشتر با مجموعه‌هایی سروکار خواهیم داشت که در آن‌ها همهٔ
!!{element}s اشیای ریاضی‌اند. چهار مجموعه از این دست آن‌قدر مهم‌اند که
نام‌های ویژه‌ای دارند:
\begin{multline*}
    \Nat = \{0, 1, 2, 3, \ldots\} \\
    \shoveright{\text{مجموعهٔ اعداد طبیعی}}\\
    \shoveleft{\Int = \{\ldots, -2, -1, 0, 1, 2, \ldots\}} \\
    \shoveright{\text{مجموعهٔ اعداد صحیح}}\\
    \shoveleft{\Rat = \Setabs{\nicefrac{m}{n}}{m, n \in \Int\text{ و }n \neq 0}}\\
    \shoveright{\text{مجموعهٔ اعداد گویا}}\\
    \shoveleft{\Real = (-\infty, \infty)}\\
    \text{مجموعهٔ اعداد حقیقی (پیوستار)}
\end{multline*}
اینها همگی مجموعه‌هایی \emph{نامتناهی} هستند؛ یعنی شمار !!{element}s
در هر یک نامتناهی است.

با گذر از هر یک از این مجموعه‌ها به مجموعهٔ بعدی، اعداد \emph{بیشتری}
به ذخیرهٔ خود می‌افزاییم. در واقع، باید روشن باشد که $\Nat \subseteq \Int
\subseteq \Rat \subseteq \Real$؛ زیرا هر عدد طبیعی یک عدد صحیح، هر
عدد صحیح یک عدد گویا، و هر عدد گویا یک عدد حقیقی است. همچنین باید روشن
باشد که $\Nat \subsetneq \Int \subsetneq
\Rat$، زیرا $-1$ عددی صحیح است اما طبیعی نیست و
$\nicefrac{1}{2}$ عددی گویاست اما صحیح نیست. این نکته کمتر آشکار است
که $\Rat \subsetneq \Real$؛ یعنی عددهای حقیقی‌ای وجود دارند
که گویا نیستند\oliflabeldef{sfr:arith:real:realline}{؛ اما در
\olref[arith][real]{realline} به این موضوع بازمی‌گردیم}{}.

گاهی از مجموعهٔ اعداد صحیح مثبت $\PosInt = \{1,
2, 3, \dots\}$ و مجموعه‌ای که تنها دو عدد طبیعی نخست را در بر دارد،
یعنی $\Bin = \{0, 1\}$، نیز استفاده خواهیم کرد.
\end{ex}

\begin{tagblock}{compsci}
\begin{ex}[رشته‌ها]
نمونهٔ جالب دیگر، مجموعهٔ $A^{*}$ از \emph{رشته‌های متناهی}
بر الفبای $A$ است: هر دنبالهٔ متناهی از عناصر~$A$ رشته‌ای بر $A$
است. \emph{رشتهٔ تهی $\Lambda$} را نیز در شمار رشته‌های بر~$A$، برای
هر الفبای~$A$، می‌آوریم. برای نمونه،
\begin{multline*}
\Bin^*
=\{\Lambda,0,1,00,01,10,11,\\
000,001,010,011,100,101,110,111,0000,\ldots\}.
\end{multline*}
اگر $x=x_{1}\ldots x_{n}\in A^{*}$ رشته‌ای متشکل از $n$
``حرف'' از $A$ باشد، می‌گوییم \emph{طول} رشته برابر~$n$ است
و می‌نویسیم $\len{x}=n$.
\end{ex}
\end{tagblock}

\begin{ex}[دنباله‌های نامتناهی]
برای هر مجموعهٔ $A$ می‌توانیم مجموعهٔ~$A^\omega$ از دنباله‌های
نامتناهی را نیز در نظر بگیریم؛ دنباله‌هایی که از !!{element}s تشکیل
می‌شوند و این اعضا به~$A$ تعلق دارند. دنبالهٔ نامتناهی $a_1a_2a_3a_4\dots$
از فهرستی نامتناهی و یک‌سویه از اشیا تشکیل شده است که هر یک
!!a{element} مجموعهٔ~$A$ است.
\end{ex}

\end{document}
